\chapter{Diskussion}

\section[Babybettchen]{Babybettchen}
\subsection[Gehäuse für die Platine]{Gehäuse für die Platine\authorlinda} \label{Platinen-Gehäuse}
Da die Platine kein wirkliches Gehäuse hat, wäre es sinnvoll ein solches zu designen und fertigen. Wir haben uns dazu bereits Gedanken gemacht und unser Ansatz wäre es mittels 3D-Druck ein Gehäuse zu fertigen. 

\par\vspace{1em}

\noindent
Der Plan war es ein Art Box zu entwerfen, in welcher die Platine und die Powerbank Platz haben. Bei dieser Box müssten an den Seitenwänden Öffnungen eingefügt werden, damit die Verkabelung der Drucksensoren zur Platine geführt werden könnten.

\noindent
Danach müsste das Gehäuse über Bohrungen mit der Holzplatte, welche sich unter der Matratze des Babybettchens befindet, verbunden werden. Dadurch könnte sich die Platine an der gleichen Stelle unter dem Babybettchen, wie zuvor befinden.

\subsection[weitere Drucksensoren]{weitere Drucksensoren\authorlinda} \label{weitere Drucksensoren}

\subsection[Herzschlag mit Drucksensoren messen]{Herzschlag mit Drucksensoren messen\authorlinda} \label{Herzschlag-Drucksensoren}

%\subsection[Wärmeerkennung]{Wärmeerkennung\authorkatrin} \label{Wärmeerkennung} maybe 

\subsection[Temperaturmessung mit Körperkontakt]{Temperaturmessung mit Körperkontakt\authorkatrin} \label{Temperaturmessung mit Körperkonkakt}
Die Messung der Körperkerntemperatur ist für eine exakte Bestimmung von Fieber oder Unterkühlung essentiell. Dafür können NTC (Negative Temparature Coefficient)
oder auch Halbleitersensoren mit integrierten ICs genutzt werden. Diese liefern oft direkt Ergrbnisse, die nicht mittels Software und berechnet werden müssen. 
Anders als bei einer Messung mittels NTC Widerständen, mit denen man mithilfe eines Spannungsteilers die Temperaturwerte aus dem jeweiligen 
Spannungswert berechnen muss.

\section[Babybettchen Code]{Babybettchen Code}
\subsection[Visualisierung der Messdaten]{Visualisierung der Messdaten\authorlinda}
Flutter vs. PWA 

\subsection[Ausbauung der PWA]{Ausbauung der PWA\authorlinda}
nicht nur localhost

%\subsection{Cloud\authorlinda}

%\section{Ki-Modell}

\section[Spannungsversorgung Testroboter]{Spannungsversorgung Testroboter\authorkatrin} \label{Spannungsversorgung Optimierung}
\section[Real Time System]{Real Time System\authorkatrin} \label{RTOS}